% =====================================================================
\chapter{Support Vector Regression}
% =====================================================================

\labinfo{11\_SVR}{svr.py}{UCI Auto MPG: 392 cars (1970--1982), 8 inputs}{Regression (fuel economy in
miles per gallon)}{RBF SVR: test RMSE 2.67 mpg, $R^2=0.887$ (least squares: RMSE 3.77, $R^2=0.776$)}

\section{Dataset}
The \textbf{Auto MPG} dataset (UCI) describes 398 car models from 1970 to 1982. The target is fuel
economy in miles per gallon (9 to 47, mean 23.4). The inputs are the number of cylinders, engine
displacement, horsepower, weight, acceleration, model year and origin (USA, Europe or Japan). Six cars
have no horsepower value, so I dropped them, leaving 392. Origin was turned into two dummy columns
(Europe, Japan) and the car name was not used. I used a random 80/20 split (314 training, 78 test cars)
and standardised the inputs.

\section{Code and experiments}
\begin{itemize}
\item \code{kernel\_matrix} gives the linear or RBF kernel.
\item \code{SVR.fit} keeps a sign vector \code{z} ($+1$ for $\alpha_i$, $-1$ for $\alpha_i^*$), a vector
      \code{pt} that says which data point each multiplier belongs to, and the running sum
      $u_k=\sum_i(\alpha_i-\alpha_i^*)K_{ik}$. Each SMO step picks the pair, computes the step \code{t},
      clips it and updates \code{u} in $O(n)$.
\item \code{predict} uses only the support vectors.
\item \code{cv\_rmse} runs 5-fold CV, standardising inside each fold.
\end{itemize}
\noindent Experiments: (1)~a grid of 48 settings, $C\in\{1,10,100,1000\}$, $\gamma\in\{0.02,0.05,0.1,0.3\}$,
$\varepsilon\in\{0.5,1,2\}$, compared by 5-fold CV RMSE; (2)~the best RBF SVR against a linear SVR and
ordinary least squares on the test set; (3)~the effect of $\varepsilon$ with $C$ and $\gamma$ fixed;
(4)~a one-input SVR (mpg against horsepower) to show the tube on real data.

\section{Results}

The best five-fold CV result was $C=100$, $\gamma=0.1$, $\varepsilon=0.5$ (CV RMSE 2.53 mpg). The next
few settings were within 0.04 mpg of it, so the choice is not delicate.

\begin{table}[H]
\centering\small
\caption{Test-set comparison (78 cars). ``In tube'' is the share of test cars predicted within $\pm\varepsilon$.}
\label{tab:svr-test}
\begin{tabular}{@{}l S[table-format=1.3] S[table-format=1.3] S[table-format=1.4] r r@{}}
\toprule
\textbf{Model} & {\textbf{MAE}} & {\textbf{RMSE}} & {\textbf{$R^2$}} & \textbf{\#SV} & \textbf{In tube}\\
\midrule
Linear regression (least squares)          & 2.749 & 3.765 & 0.7759 & -- & --\\
Linear SVR ($C=10$, $\varepsilon=1$)       & 2.798 & 4.029 & 0.7434 & 221 & 24.4\%\\
RBF SVR ($C=100$, $\gamma=0.1$, $\varepsilon=0.5$) & \bfseries 1.936 & \bfseries 2.672 & \bfseries 0.8871 & 241 & 29.5\%\\
\bottomrule
\end{tabular}
\end{table}

The RBF SVR cuts the test RMSE by 29\% compared with least squares (3.77 to 2.67 mpg) and raises $R^2$
from 0.78 to 0.89. Fuel economy depends on weight and horsepower in a curved way (it falls quickly at
first and then levels off), which a straight line cannot capture. The linear SVR is not better than least
squares; it is still a straight line, just fitted with a different loss. So the improvement comes from
the kernel, not from the loss function.

\begin{table}[H]
\centering\small
\caption{Effect of the tube width $\varepsilon$ ($C=100$, $\gamma=0.1$).}
\label{tab:svr-eps}
\begin{tabular}{@{}l c c c c c c c c@{}}
\toprule
$\varepsilon$ (mpg) & 0.1 & 0.25 & 0.5 & 1 & 2 & 3 & 4 & 6\\
\midrule
Support vectors (of 314) & 300 & 276 & 241 & 193 & 111 & 67 & 44 & 24\\
Test RMSE (mpg) & 2.681 & 2.657 & 2.672 & 2.673 & 2.820 & 3.038 & 3.283 & 3.809\\
\bottomrule
\end{tabular}
\end{table}

The $\varepsilon$ study shows the sparsity of SVR clearly. With $\varepsilon=1$~mpg the model uses 193
support vectors and is as accurate as with $\varepsilon=0.1$ (300 support vectors). With
$\varepsilon=4$~mpg only 44 cars define the whole function, at the cost of about 0.6 mpg of extra RMSE.

\begin{figure}[H]
\centering
\includegraphics[width=0.48\linewidth]{\figpath{11\_SVR}{epsilon_tube_1d.pdf}}\hfill
\includegraphics[width=0.48\linewidth]{\figpath{11\_SVR}{pred_vs_actual.pdf}}
\caption{The SVR epsilon tube and predicted-versus-actual test values.}
\end{figure}

\section{Source code}
\lstinputlisting[caption={\texttt{11\_SVR/svr.py}}]{11_SVR/svr.py}
\clearpage
